\section{DP 四问法}\label{dp-ux56dbux95eeux6cd5}

\begin{enumerate}
\def\labelenumi{\arabic{enumi}.}
\tightlist
\item
  状态表示什么最小信息？
\item
  最后一步/最后一个决策是什么？
\item
  转移是否覆盖且不重复所有情况？
\item
  初值、遍历顺序、答案位置是什么？
\end{enumerate}

先写朴素正确版本，再按状态数量、转移数量、内存逐层优化。\passthrough{\lstinline!dp!} 数组含义要写在代码注释里。

\section{线性 DP}\label{ux7ebfux6027-dp}

\begin{itemize}
\tightlist
\item
  LIS：\passthrough{\lstinline!O(n log n)!} 维护每个长度的最小结尾。
\item
  严格递增用 \passthrough{\lstinline!lower\_bound!}，非降用 \passthrough{\lstinline!upper\_bound!}。
\item
  最长公共子序列：\passthrough{\lstinline!dp[i][j]!} 看前缀；滚动到一维时保存左上角旧值。
\item
  最大子段和：\passthrough{\lstinline!bestEnding = max(a[i], bestEnding+a[i])!}，也可用前缀最小值。
\item
  计数 DP：同一状态的方案相加，注意顺序是否导致重复计数。
\end{itemize}

\begin{lstlisting}[language={C++}]
// 变量：tail=LIS 各长度对应的最小结尾值数组 tail。
vector<int> tail;
// 变量：x=从 a 枚举的元素值 x。
for (int x : a) {
    auto it = lower_bound(tail.begin(), tail.end(), x); // 非降改 upper_bound
    if (it == tail.end()) tail.push_back(x);
    else *it = x;
}
// 变量：lis=当前最长递增子序列长度 lis。
int lis = tail.size();

// 变量：bestEnding=必须以当前位置结尾的精确最大子段和；maxSubarray=扫描至今的精确最大子段和。
__int128 bestEnding = a[0], maxSubarray = a[0];
// 变量：i=最大子段和扫描中当前作为结尾的数组位置 i。
for (int i = 1; i < n; ++i) {
    bestEnding = max<__int128>(a[i], bestEnding + a[i]);
    maxSubarray = max(maxSubarray, bestEnding);
}
\end{lstlisting}

\begin{lstlisting}[language={C++}]
// 变量：lcs=lcs[i][j] 保存 a 前 i 个字符与 b 前 j 个字符的最长公共子序列长度。
vector<vector<int>> lcs(n + 1, vector<int>(m + 1));
// 变量：i=LCS 状态 lcs[i][j] 对应的字符串 a 前缀长度 i。
for (int i = 1; i <= n; ++i)
    // 变量：j=LCS 状态 lcs[i][j] 对应的字符串 b 前缀长度 j。
    for (int j = 1; j <= m; ++j)
        lcs[i][j] = (a[i - 1] == b[j - 1])
            ? lcs[i - 1][j - 1] + 1
            : max(lcs[i - 1][j], lcs[i][j - 1]);
\end{lstlisting}

\section{背包}\label{ux80ccux5305}

\begin{longtable}[]{@{}ll@{}}
\toprule\noalign{}
类型 & 状态/遍历 \\
\midrule\noalign{}
\endhead
\bottomrule\noalign{}
\endlastfoot
0/1 背包 & \passthrough{\lstinline!dp[j]=max(dp[j],dp[j-w]+v)!}，\passthrough{\lstinline!j!} 从大到小 \\
完全背包 & 同式，\passthrough{\lstinline!j!} 从小到大 \\
多重背包 & 二进制拆分或单调队列优化 \\
分组背包 & 每组只能选一个，组内倒序并用上一层状态 \\
依赖背包 & 树形 DP/附件分组 \\
计数背包 & 把 \passthrough{\lstinline!max!} 换成加法，模数下处理重复顺序 \\
\end{longtable}

容量 \passthrough{\lstinline!W≤1e5!} 时常用 \passthrough{\lstinline!O(nW)!}；物品价值小则把``最小重量''作为状态；\passthrough{\lstinline!n≤40!} 可折半搜索。

\begin{lstlisting}[language={C++}]
// 接口：zero_one_knapsack(items, W)：返回每个容量上限下的 0/1 背包最大价值。
// 参数：W>=0；每件物品重量为正且价值为 int。正重量保证可选件数不超过 W，long long 足以容纳最优价值。
vector<long long> zero_one_knapsack(const vector<pair<int, int>>& items, int W) {
    // 变量：dp=容量不超过下标时的 0/1 背包最大价值。
    vector<long long> dp(W + 1, 0);
    // 变量：w=当前物品重量；v=当前物品价值。
    for (auto [w, v] : items)
        // 变量：j=倒序枚举的背包容量。
        for (int j = W; j >= w; --j)
            dp[j] = max(dp[j], dp[j - w] + v);
    return dp;
}
// 接口：complete_knapsack(items, W)：返回每个容量上限下的完全背包最大价值。
// 参数：W>=0；每件物品重量为正且价值为 int。最多选择 W 件，long long 足以容纳最优价值。
vector<long long> complete_knapsack(const vector<pair<int, int>>& items, int W) {
    // 变量：dp=容量不超过下标时的完全背包最大价值。
    vector<long long> dp(W + 1, 0);
    // 变量：w=当前物品重量；v=当前物品价值。
    for (auto [w, v] : items)
        // 变量：j=正序枚举的背包容量。
        for (int j = w; j <= W; ++j)
            dp[j] = max(dp[j], dp[j - w] + v);
    return dp;
}
\end{lstlisting}

\section{区间 DP}\label{ux533aux95f4-dp}

按区间长度递增，枚举分割点：\passthrough{\lstinline!dp[l][r] = min/max(dp[l][k]+dp[k+1][r]+cost)!}。典型：石子合并、矩阵链乘、括号匹配、回文分割。若代价满足四边形不等式/决策单调性，可用四边形不等式优化，但先验证条件。

\begin{lstlisting}[language={C++}]
// 变量：len=区间 DP 当前枚举的区间长度 len。
for (int len = 2; len <= n; ++len)
    // 变量：l=当前区间左端点 l。
    for (int l = 1; l + len - 1 <= n; ++l) {
        // 变量：r=当前区间右端点 r。
        int r = l + len - 1;
        dp[l][r] = INF;
        // 变量：k=当前循环枚举的整数下标 k。
        for (int k = l; k < r; ++k)
            dp[l][r] = min(dp[l][r], dp[l][k] + dp[k + 1][r] + cost(l, r));
    }
\end{lstlisting}

\section{树形 DP与换根 DP}\label{ux6811ux5f62-dpux4e0eux6362ux6839-dp}

树上没有环，后序即可保证子树状态已知。独立集：\passthrough{\lstinline!f[u][0]=Σmax(f[v][0],f[v][1])!}，\passthrough{\lstinline!f[u][1]=w[u]+Σf[v][0]!}。换根时先从任意根求子树贡献，再将父亲的答案传给儿子，常见于``每个点到其他点距离和''。

\begin{lstlisting}[language={C++}]
// 接口：reroot1(u, p)：第一次 DFS，计算每个子树的向下 DP；不返回值。
// 参数：当前顶点编号 u；当前节点 u 的父节点 p。
void reroot1(int u, int p) {
    sz[u] = 1; down[u] = 0;
    // 变量：v=从 tree[u] 枚举的相邻顶点 v。
    for (int v : tree[u]) if (v != p)
        reroot1(v, u), sz[u] += sz[v], down[u] += down[v] + sz[v];
}
// 接口：reroot2(u, p)：第二次 DFS，把父侧贡献传给儿子并得到全树答案；不返回值。
// 参数：当前顶点编号 u；当前节点 u 的父节点 p。
void reroot2(int u, int p) {
    // 变量：v=从 tree[u] 枚举的相邻顶点 v。
    for (int v : tree[u]) if (v != p) {
        ans[v] = ans[u] - sz[v] + (n - sz[v]);
        reroot2(v, u);
    }
}
\end{lstlisting}

\section{有向无环图动态规划与路径计数（Directed Acyclic Graph Dynamic Programming）}\label{dag-dpux4e0eux8defux5f84ux8ba1ux6570}

拓扑序内转移；最长路用 \passthrough{\lstinline!-INF!} 初始化，最短路用 \passthrough{\lstinline!INF!}。若有多个源，把所有源初始化；要恢复路径保存 \passthrough{\lstinline!pre!}。在模数下统计方案时不要把不可达状态的 \passthrough{\lstinline!0!} 与真实方案混淆。

\begin{lstlisting}[language={C++}]
// 变量：u=从容器 topo 枚举的元素 u。
for (int u : topo) if (dp[u] != NEG)
    // 变量：v=当前边的终点 v；w=当前边权 w。
    for (auto [v, w] : dag[u]) {
        // 变量：candidate=经 u 到达 v 的候选最长路长度 candidate。
        long long candidate = dp[u] + w;
        if (candidate > dp[v]) {
            dp[v] = candidate;
            ways[v] = ways[u];
        } else if (candidate == dp[v]) {
            ways[v] = (ways[v] + ways[u]) % MOD;
        }
    }
\end{lstlisting}

\section{状态压缩动态规划}\label{ux72b6ux538b-dp}

\passthrough{\lstinline!mask!} 表示已选集合，转移 \passthrough{\lstinline!mask | (1<<i)!}。\passthrough{\lstinline!n≤20!} 的旅行商、集合覆盖、配对、棋盘放置都适用。优化：预处理 \passthrough{\lstinline!popcount!}、子集和、可行转移；枚举子集总复杂度 \passthrough{\lstinline!O(3\^n)!}。

\begin{lstlisting}[language={C++}]
// 变量：dp=dp[mask][u] 表示从顶点 0 出发、恰好访问 mask 中顶点且停在 u 的最短路径代价。
vector<vector<long long>> dp(1 << n, vector<long long>(n, INF));
dp[1][0] = 0;
// 数值域：边权非负；INF 是不参与转移的上界，所有需要区分的有限路径代价必须严格小于 INF。
// 变量：mask=用二进制位编码的当前状态 mask。
for (int mask = 1; mask < (1 << n); ++mask)
    // 变量：u=状态压缩 DP 当前路径末端的顶点 u。
    for (int u = 0; u < n; ++u) if (mask >> u & 1)
        // 变量：v=状态压缩 DP 当前尝试加入集合的顶点 v。
        for (int v = 0; v < n; ++v) if (!(mask >> v & 1) && dp[mask][u] < INF) {
            // 变量：candidate=宽整数计算的候选路径代价；达到 INF 后按本接口的上界处理。
            __int128 candidate = (__int128)dp[mask][u] + w[u][v];
            if (candidate < dp[mask | (1 << v)][v])
                dp[mask | (1 << v)][v] = (long long)candidate;
        }
\end{lstlisting}

\section{数位动态规划（Digit Dynamic Programming）}\label{ux6570ux4f4d-dp}

按高位到低位处理，状态通常为 \passthrough{\lstinline!(pos, tight, started, remainder/计数)!}。\passthrough{\lstinline!tight!} 表示前缀是否贴合上界，\passthrough{\lstinline!started!} 区分前导零。求 \passthrough{\lstinline![L,R]!} 用 \passthrough{\lstinline!F(R)-F(L-1)!}。先写``不限制上界''的转移，再加入 tight。

\begin{lstlisting}[language={C++}]
long long memo[20][2][2][11]; // -1=未计算；末维 rem 是模 11 余数 [0,10]。
string digits; // 上界的十进制表示；pos 从 0 扫到 digits.size()-1。
// 接口：solve_digit(pos, tight, started, rem)：返回给定数位 DP 状态的合法后缀方案数；返回类型 long long。
// 参数：当前处理到的十进制数位下标 pos；前缀是否仍贴住上界 tight；是否已经放置非前导零数字 started；当前余数状态 rem。
long long solve_digit(int pos, int tight, int started, int rem) {
    if (pos == (int)digits.size()) return started && rem == 0;
    // 变量：res=当前 DP 状态对应的记忆化答案引用 res。
    long long &res = memo[pos][tight][started][rem];
    if (!tight && res != -1) return res;
    res = 0;
    // 变量：lim=数位 pos 可枚举数字的上界 lim；tight 为真时取 digits[pos]，否则为 9。
    int lim = tight ? digits[pos] - '0' : 9;
    // 变量：d=数位 DP 当前枚举的十进制数字 d。
    for (int d = 0; d <= lim; ++d)
        res += solve_digit(pos + 1, tight && d == lim, started || d,
                           (rem * 10 + d) % mod);
    return res;
}
\end{lstlisting}

\section{概率动态规划与期望动态规划（Probability and Expected-Value Dynamic Programming）}\label{ux6982ux7387-dpux4e0eux671fux671b-dp}

状态转移可能含当前状态自身：整理成 \passthrough{\lstinline!E = (1 + other)/(1-p\_self)!}。马尔可夫链有限状态可高斯消元；若转移无环则按拓扑/记忆化直接递推。概率取模时乘逆元，浮点题要保留足够精度。

\begin{lstlisting}[language={C++}]
// 无环状态图上的期望：先按题意保证 next[u] 都比 u 更靠后
// 变量：E=E[u] 表示从无环状态图的状态 u 出发直到终止所需的期望步数。
vector<double> E(n, 0.0);
// 变量：u=按逆拓扑序计算期望的当前状态 u。
for (int u = n - 1; u >= 0; --u) {
    if (transitions[u].empty()) continue; // 终止态没有后续操作，剩余代价为 0。
    E[u] = 1.0;
    // 变量：v=转移到的后继状态编号 v；p=采用该转移的概率 p。
    for (auto [v, p] : transitions[u]) E[u] += p * E[v];
}
\end{lstlisting}

\section{记忆化搜索}\label{ux8bb0ux5fc6ux5316ux641cux7d22}

当状态图是 DAG 或状态数量可控时，用 DFS + cache 比显式排序更自然。必须保证每个状态的所有后继已计算，含环问题不能直接记忆化递归。

\begin{lstlisting}[language={C++}]
int memo[MAXN]; // -1=状态尚未计算；否则为 solve(u) 的缓存答案。
// 接口：solve(u)：记忆化返回从 DAG 状态 u 出发可取得的最大后继收益；返回类型 int。
// 参数：顶点编号 u。
int solve(int u) {
    // 变量：res=当前 DP 状态对应的记忆化答案引用 res。
    int &res = memo[u];
    if (res != -1) return res;
    res = 0;
    // 变量：v=从 next[u] 枚举的后继状态 v。
    for (int v : next[u]) res = max(res, solve(v) + gain[u][v]);
    return res;
}
\end{lstlisting}

\section{动态规划优化（Dynamic Programming Optimization）}\label{dp-ux4f18ux5316}

\subsection{滚动数组}\label{ux6edaux52a8ux6570ux7ec4}

只依赖上一层时把二维压一维；确认本轮不会读到已更新的状态，遍历方向是关键。

\begin{lstlisting}[language={C++}]
// 变量：prev=滚动数组中保存的上一 DP 层 prev；cur=滚动数组中正在计算的当前 DP 层 cur。
vector<long long> prev(m + 1), cur(m + 1);
// 变量：i=滚动数组中正在计算的当前 DP 层编号 i。
for (int i = 1; i <= n; ++i) {
    cur[0] = 0;
    // 变量：j=滚动数组中正在计算的当前 DP 状态下标 j。
    for (int j = 1; j <= m; ++j)
        cur[j] = transition(prev[j], cur[j - 1], a[i], b[j]);
    swap(prev, cur);
}
\end{lstlisting}

\subsection{单调队列优化}\label{ux5355ux8c03ux961fux5217ux4f18ux5316}

形如 \passthrough{\lstinline!dp[i]=base(i)+min/max(dp[j]+g(j))!} 且 \passthrough{\lstinline!j!} 的合法范围滑动时，维护候选队列。若 \passthrough{\lstinline!g!} 可分离且窗口固定，复杂度从 \passthrough{\lstinline!O(nk)!} 降到 \passthrough{\lstinline!O(n)!}。

\begin{lstlisting}[language={C++}]
// 接口：window_min_dp(base, k)：在 __int128 中精确计算 dp[i]=base[i]+min(dp[j])，其中 i-k<=j<i；dp[0]=base[0]。
// 参数：每个位置的附加代价 base；正窗口宽度 k。
vector<__int128> window_min_dp(const vector<long long>& base, int k) {
    // 变量：n=DP 状态总数。
    int n = (int)base.size();
    if (n == 0) return {};
    // 变量：dp=窗口最小值优化后的状态数组。
    vector<__int128> dp(n);
    // 变量：q=按 dp 值递增维护的候选下标队列。
    deque<int> q;
    dp[0] = base[0]; q.push_back(0);
    // 变量：i=当前计算的 DP 状态下标。
    for (int i = 1; i < n; ++i) {
        while (!q.empty() && q.front() < i - k) q.pop_front();
        dp[i] = base[i] + dp[q.front()];
        while (!q.empty() && dp[q.back()] >= dp[i]) q.pop_back();
        q.push_back(i);
    }
    return dp;
}
\end{lstlisting}

\subsection{斜率优化（Convex Hull Trick，CHT）}\label{ux659cux7387ux4f18ux5316ux51f8ux5305-trick}

转移能化为 \passthrough{\lstinline!dp[i]=x\_i*k\_j+b\_j!}，把每个 \passthrough{\lstinline!j!} 看作直线，按斜率/查询单调性维护凸包。斜率不单调时用 Li Chao 树；确认除法比较不溢出，必要时用 \passthrough{\lstinline!\_\_int128!}。

\begin{lstlisting}[language={C++}]
struct Line {
    long long k, b; // 候选直线 y=k*x+b 的斜率与截距。
    // 接口：get(x)：在 __int128 中精确返回该直线在横坐标 x 处的函数值 k*x+b。
    // 参数：任意 long long 查询横坐标 x。
    __int128 get(long long x) const { return (__int128)k * x + b; }
};
// 变量：hull=按有效顺序维护的凸包候选队列 hull。
deque<Line> hull;
// 接口：bad(a, b, c)：判断中间直线 b 是否因交点顺序而永远不会成为最优；可删除 b 时返回 true。
// 参数：按斜率顺序相邻的候选直线 a、b、c；要求两个交叉乘积可由 signed __int128 表示，差值会先提升再相减。
bool bad(Line a, Line b, Line c) {
    return ((__int128)b.b - a.b) * ((__int128)b.k - c.k)
         >= ((__int128)c.b - b.b) * ((__int128)a.k - b.k);
}
// 接口：add_line(x)：按单调斜率加入候选直线，并从队尾删除永不最优的直线；会修改当前结构；多测时重新构造或显式清空成员。
// 参数：待插入的候选直线 x。
void add_line(Line x) {
    while (hull.size() >= 2 && bad(hull[hull.size() - 2], hull.back(), x)) hull.pop_back();
    hull.push_back(x);
}
\end{lstlisting}

\subsection{分治优化}\label{ux5206ux6cbbux4f18ux5316}

若最优决策点单调，\passthrough{\lstinline!dp[i][j]!} 在分治递归中只枚举决策区间的一部分，常见 \passthrough{\lstinline!O(kn log n)!}。先证明四边形不等式或决策单调性。

\begin{lstlisting}[language={C++}]
// 接口：solve(l, r, optl, optr)：在决策范围 [optl,optr] 内计算 DP 下标区间 [l,r] 并递归缩小最优决策范围；无返回值。
// 参数：闭区间左端点 l；闭区间右端点 r；真实最优决策允许范围左端 optl；真实最优决策允许范围右端 optr。
void solve(int l, int r, int optl, int optr) {
    if (l > r) return;
    // 变量：mid=当前区间中点 mid；best=目前找到的最优值 best。
    int mid = (l + r) / 2, best = optl;
    // 变量：k=当前 DP 状态枚举的候选决策下标 k。
    for (int k = optl; k <= min(mid, optr); ++k)
        if (cost(k, mid) < cost(best, mid)) best = k;
    dp[mid] = cost(best, mid);
    solve(l, mid - 1, optl, best);
    solve(mid + 1, r, best, optr);
}
\end{lstlisting}

\section{组合游戏动态规划（Combinatorial Game Dynamic Programming）}\label{ux7ec4ux5408ux6e38ux620f-dp}

无环游戏状态标记胜负：存在一个后继为必败则当前必胜；所有后继必胜则当前必败。将多个独立游戏的 Grundy 值 xor。棋盘/取石子题先找周期或 SG。

\begin{lstlisting}[language={C++}]
// 接口：grundy(x)：返回状态参数对应的 Sprague-Grundy 值并记忆化；返回类型 int。
// 参数：待求 SG 值的游戏状态 x。
int grundy(int x) {
    if (memo[x] != -1) return memo[x];
    // 变量：seen=记录当前状态各后继 SG 值是否出现的布尔表 seen。
    bool seen[MAXG] = {};
    // 变量：y=状态 x 可到达的后继状态 y。
    for (int y : moves[x]) seen[grundy(y)] = true;
    // 变量：g=当前有向图的邻接表 g。
    int g = 0; while (seen[g]) ++g;
    return memo[x] = g;
}
\end{lstlisting}

\section{常见 DP 误区}\label{ux5e38ux89c1-dp-ux8befux533a}

\begin{itemize}
\tightlist
\item
  用``最后一个元素''定义状态，却漏掉了影响未来的最小信息。
\item
  把排列计数当组合计数，循环顺序写反。
\item
  区间 DP 长度顺序错误。
\item
  负权最大化仍用 \passthrough{\lstinline!0!} 初始化，导致不可达状态污染答案。
\item
  取模减法忘记加模；概率 DP 把浮点误差当精确相等。
\end{itemize}
